\subsection{An explicit characteristic source for the character}

\begin{result}{A.4.2}{Proved}{A finite characteristic modular-source detector for every multiplier}
Put
\[
 n_*=2^{51942}3^{14136},\qquad
 m=576n_*^{97}=2^{5038380}3^{1371194}.
\]
In the full profinite modular group
$\widehat B=\langle s,r\mid s^2=r^3=1\rangle^\wedge$, define
\[
 K_m=\bigcap_{\theta\in
       \operatorname{Hom}_{\rm cts}(\widehat B,\mathfrak S_m)}
                      \ker\theta .
 \tag{A.H1}
\]
This is an explicitly specified open characteristic subgroup.
For the genuine normalized gentle action on $\widehat B$, the character
$\Delta$ factors through the finite action on $\widehat B/K_m$.
The quotient has order at most $(m!)^{(m!)^2}$.
Before taking the characteristic core, the smaller gentle-stable
quotient $\widehat B/(N'N^{n_*})$ has order $m$, where
$N=\ker(\widehat B\to\operatorname{SL}_2(\mathbb Z/12)/\{\pm I\})$.
This proves the source-detection part of task A.4; the five-point
enumeration in A.4 remains a separate unresolved calculation.
\end{result}
\begin{proof}
We include the maps and the full-multiplier argument. The genuine
three-strand action provided by the refereed R1 normalization is
\[
 u\longmapsto u^\lambda,\qquad
 w\longmapsto f(u^2,w^2)^{-1}w^\lambda f(u^2,w^2).
\]
Let $q(u)=U=\left(\begin{smallmatrix}1&1\\0&1\end{smallmatrix}\right)$
and $q(w)=L=\left(\begin{smallmatrix}1&0\\-1&1\end{smallmatrix}\right)$.
For each prime $p$, write the image of the second generator as
$B=\left(\begin{smallmatrix}a&b\\c&2-a\end{smallmatrix}\right)$
over $\mathbb Z_p$. The braid relation with $A=U^\lambda$ gives
$c(1+\lambda c)=0$. The case $c=0$ forces $a=1,b=\lambda$
by the determinant and upper-right braid equations, hence $B=A$.
The action is an automorphism and the original image generated by
$U,L$ is nonabelian, so this case is excluded. Thus
$c=-\lambda^{-1}$. With $w_0=1-a$, the same equations give
\[
 B=\begin{pmatrix}1-w_0&\lambda w_0^2\\
                  -\lambda^{-1}&1+w_0\end{pmatrix}
   =VLV^{-1},\qquad
 V=\begin{pmatrix}\lambda&\lambda w_0\\0&1\end{pmatrix}.
\]
Also $VUV^{-1}=U^\lambda$. The simultaneous centralizer of $U,L$
consists of scalars: commuting with $U$ forces the lower-left
entry to vanish and the two diagonal entries to agree, and then
commuting with $L$ kills the upper-right entry. Hence $V$ is unique
with lower-right entry one. Taking the product over all primes
gives a continuous adelic covariance homomorphism, with the actual
composition order inherited from composition of the genuine actions.
Reduction modulo twelve proves that $N$ is preserved by every
gentle automorphism. No characteristicity of $N$ under arbitrary
automorphisms is needed in this step.

Write $Q=\operatorname{SL}_2(\mathbb Z/12)/\{\pm I\}$,
$H_2=\langle q(s)\rangle$, $H_3=\langle q(r)\rangle$.
The action of the discrete congruence kernel on the Bass--Serre
tree is free. Its quotient graph has edges $Q$ and vertices
$Q/H_2\amalg Q/H_3$, so its first homology has rank
$576-288-192+1=97$. The full induced topology on this
finite-index free subgroup identifies its completed homology with
$M=N^{\rm ab}=\widehat{\mathbb Z}^{97}$. The integral graph
complex and its profinite completion give the exact $Q$-module
sequence
\[
 0\to M\to \widehat{\mathbb Z}[Q]
 \to \widehat{\mathbb Z}[Q/H_2]\oplus
       \widehat{\mathbb Z}[Q/H_3]
 \to\widehat{\mathbb Z}\to0.                              \tag{A.H2}
\]
Its completion is exact because the graph homology lattices are
free. Let $L_0$ be the kernel of the last map. The invariant
augmentation is $(a,b)\mapsto288a+192b$, and
$L_0^Q=\widehat{\mathbb Z}(2,-3)$. The invariant edge norm
maps to this primitive generator. Shapiro's lemma on the two
short exact sequences in (A.H2) therefore yields
\[
 H^1(Q,M)=0,\qquad H^2(Q,M)=\mathbb Z/96.                  \tag{A.H3}
\]
Explicitly, $H^{>0}(Q,\widehat{\mathbb Z}[Q])=0$,
$H^1(Q,\widehat{\mathbb Z}[Q/H_i])=
\operatorname{Hom}(H_i,\widehat{\mathbb Z})=0$, and the
remaining cokernel is that of $(288,192)$.

For $g\ne1$, the trace on $M$ is
$1-\#(Q/H_2)^g-\#(Q/H_3)^g$. A nonidentity element cannot
fix a coset of both types because its order would divide both
two and three. Thus
$96\le97-\operatorname{tr}(g\mid M)\le384<576$.
The action of $Q$ on $M/576M$ is faithful. The program
\texttt{check\_homology\_detector.py} constructs the graph,
a spanning tree and all $97$ integral cycles. It independently
checks both deck matrices, all $576$ traces, and the primitive
cusp cycle. The exact trace distribution is
\[
 \operatorname{tr}=-7,-5,1,97
 \quad\text{with multiplicities }36,64,475,1.             \tag{A.H4}
\]

Let $h_{12,n}$ denote the induced genuine action on $M/nM$.
First take a multiplier-one element $\Phi$ with
$h_{12,576}(\Phi)=1$. Equivariance and the just-proved
faithfulness force $\Phi_Q=1$. In
$\Gamma_n=\widehat B/(N'N^n)$, the difference
$\Phi(b)b^{-1}$ for a lift $b$ of $g\in Q$ is a well-defined
cocycle in $Z^1(Q,M/576M)$. The coefficient sequence and
(A.H3) identify $H^1(Q,M/nM)$ with $H^2(Q,M)[n]$.
Reduction from $576$ to $6$ is multiplication by $96$ under
this identification: the left vertical map in the two
coefficient sequences is multiplication by $96$, and the
middle one is the identity. The reduction class vanishes.
Consequently the action on $\Gamma_6$ is conjugation by an
element of $M/6M$.

To read the frames, use the actual image $E_0$ of $\widehat B$
in
\[
 \bigl(\mathbb F_2[Q/H_2]\oplus\mathbb F_3[Q/H_3]\bigr)
                         \rtimes Q,
\]
where
$s\mapsto(e_{H_2},0,q(s))$ and
$r\mapsto(0,e_{H_3},q(r))$. These images have orders two and
three, respectively, so they define the stated continuous
homomorphism. Its restriction to $N$ is abelian of exponent
dividing six, and therefore factors through $\Gamma_6$.
The projection of $C_{E_0}(s)$ is contained in $H_2$:
the first-coordinate centralizer equation, read at the
$s$-fixed coset $H_2$, forces $qH_2=H_2$. In detail,
$(1-s)v=e_{H_2}-e_{qH_2}$ has left coefficient zero at
$H_2$, whereas the right coefficient there is one unless
$qH_2=H_2$. The same calculation in characteristic three
proves that the projection of $C_{E_0}(r)$ is contained in
$H_3$.

The positive pure frames satisfy
$\Phi(s)=\beta s\beta^{-1}$ and $\Phi(r)=ArA^{-1}$.
Comparing these with the common inner action on $E_0$
gives $q(\beta)\in H_2$ and $q(A)\in H_3$.
Purity and the quotient to $\mathfrak S_3$ make both images
one. The R1 signed normalization is $I$ modulo four,
excluding the other lift of the projective identity.
The frame abelianization gives $\rho=0$ modulo three,
so $U^{8\rho}=I$ modulo twelve. Thus the two signed
corrected frames are $T=W=I$, and $\Delta(\Phi)=0$.
This proves the multiplier-one detection assertion.

We next remove the multiplier restriction by an actual arithmetic
normalization. The quotient $F/(N'N^{576})$ has order with
prime divisors only two and three. The genuine tangential action
is the restriction of the displayed braid action to
$x=u^2,y=w^2$. For a rational prime $p\ge5$, the smooth
pair $(\mathbb P^1_{\mathbb Z_p},\{0,1,\infty\})$ and the
unit tangent $\overrightarrow{01}$ satisfy the genuine
prime-to-$p$ specialization comparison
\cite[Lemma 8.7]{Petrov2023}. Its action factors through
residue Galois, so inertia acts trivially on this quotient
and hence on $M/576M$. Therefore the fixed field $L$ of
\[
 K=\ker(h_{12,576}|_{G_{\mathbb Q}})
\]
is finite Galois, unramified outside two and three, with
degree dividing $D=|\operatorname{GL}_{97}(\mathbb Z/576)|$.

Let $\Lambda=\chi(K)\subset\widehat{\mathbb Z}^{\times}$.
The restriction Galois correspondence gives
$[\widehat{\mathbb Z}^{\times}:\Lambda]
 =[L\cap\mathbb Q^{\rm cyc}:\mathbb Q]\mid D$.
Cyclotomic inertia at $p$ is exactly the $p$-factor of
$\prod_p\mathbb Z_p^\times$. Thus $\Lambda$ contains all
factors at $p\ge5$, and closedness gives their product.
In the remaining finite quotient, the images of
$1+4\mathbb Z_2$ and $1+3\mathbb Z_3$ have orders dividing
$2^{v_2(D)}$ and $3^{v_3(D)}$. The principal-unit power
identities give containment of
$1+2^{v_2(D)+2}\mathbb Z_2$ and
$1+3^{v_3(D)+1}\mathbb Z_3$ in $\Lambda$.
The exact general-linear order formula gives
\[
 \begin{split}
 v_2(D)&=5(97)^2+97\cdot96/2+
           \sum_{i=1}^{97}v_2(3^i-1)=51940,\\
 v_3(D)&=(97)^2+97\cdot96/2+
           \sum_{i=1}^{97}v_3(2^i-1)=14135.
 \end{split}                                             \tag{A.H5}
\]
The two sums are $239$ and $70$; the program verifies them
by exact integer arithmetic. It follows that the principal
congruence kernel of modulus $n_*$ lies in $\Lambda$.

The class of $\tau^{12}$ in $M$ is primitive. Its graph
cycle is embedded and has coefficients $1$ and $-1$; the
program also computes its coordinates in the chosen integral
basis. Since $\Phi(\tau)=\tau^\lambda$, an element in
$\ker h_{12,n_*}$ has $\lambda\equiv1\pmod{n_*}$.
Choose $\sigma\in K$ with exactly that full cyclotomic
multiplier, using the membership $\lambda\in\Lambda$.
Then $\Phi\sigma^{-1}$ has multiplier one and trivial
$h_{12,576}$-action. The preceding paragraph and R1 give
$\Delta(\Phi\sigma^{-1})=\Delta(\sigma)=0$.
This proves detection by $h_{12,n_*}$ for all multipliers.

Finally $\Gamma_{n_*}$ has order $576n_*^{97}=m$.
Its regular permutation representation is one of the
homomorphisms in (A.H1), so
$K_m\subset N'N^{n_*}$. The intersection is finite, because
a continuous homomorphism is determined by its two generator
images; there are at most $(m!)^2$ such pairs. Precomposition
by any continuous automorphism permutes these homomorphisms,
which proves characteristicity. The diagonal product embeds
$\widehat B/K_m$ into at most $(m!)^2$ copies of
$\mathfrak S_m$, giving the asserted bound.
An automorphism acting trivially on this characteristic
quotient acts trivially on $\Gamma_{n_*}$ and on $M/n_*M$.
The kernel containment just proved is exactly the desired
factorization of $\Delta$.
\end{proof}
